\section{Background and Motivation}

\subsection{Three Roads to Eternal Inflation}

% In order that matter and galaxies could to have come to populate our observable universe, inflation must have come to an end within our Hubble volume.  The timing of that transition may in general vary across large distance scales with the local behavior of the inflaton fields, leading to the partition of space into non-inflating {\it pocket universes} embedded in an inflating background.  If the inflating volume consistently reproduces faster than these pocket universes can form, then the cycle will continue without indefinitely.  This scenario is known as {\it eternal inflation}.

Eternal inflation may manifest in any of three discernable modes: stochastic, topological, and false-vacuum.  More than one mode may manifest in a given inflation model, but only one of the associated criteria need be be satisfied to all but ensure ever-lasting inflation.  As such, those wishing to mount a defense against the creeping inevitability of eternal inflation must be guarded on three sides.

\textbf{Stochastic eternal inflation} occurs when quantum fluctuations of the inflaton field dominate over its classical evolution according to the slow roll equation of motion.
Over the passage of a Hubble time $H^{-1}$, the field expectation value $\phi$ decreases by an amount $\Delta \phi = \lvert\dot\phi\rvert H^{-1} = V_{,\phi} / 3H^2 $ as it classically rolls.  During the same interval, quantum fluctuations with amplitude $\delta\phi \approx H/2\pi$ may drive the field up or down the potential slope relative to its classical trajectory.  
If in that time, the probability for the field within an initially Hubble-sized region to move up the slope is greater than the ratio of initial to final volumes ($1/e^3$), then at least one final Hubble volume is likely to continue inflating after a time $H^{-1}$.  This occurs when $ \Delta \phi \lesssim \delta \phi \, \sqrt{2} \, \erfc^{-1} (2e^{-3}) $ or
\begin{equation}
% \delta\phi \gtrsim 0.61 \, \Delta\phi 
V^{3/2} \geq 6.64 \, V_{,\phi} \label{seic}
\end{equation}  
We refer to this as the {\it stochastic eternal inflation criterion}.

\textbf{Topological eternal inflation} may manifest in cases when spatially inhomogeneous field solutions are divisible into (approximately) inequivalent homotopy classes.  
Consider for example a symmetric potential with two degenerate vacua at $\phi = \pm \phi_0$, separated by a local maximum at $\phi = 0$.  
If in neighboring regions the field has settled into different vacua, then by continuity the field must reach the maximum of the potential somewhere in between.  
This configuration cannot classically evolve into a true vacuum solution, as that would require the energy to increase as the field in one region traverses the central maximum.
If the space occupying the middle ground is undergoing inflationary expansion at a rate sufficient to compensate for the recession of its boundaries, then we have {topological eternal inflation}.

In \textbf{false-vacuum eternal inflation}, the field becomes trapped near a local minimum of the potential, corresponding to a positive vacuum energy.  Classically, the field would be stuck there indefinitely.  However, quantum effects allow it to ``escape'' by tunneling out of this false vacuum to continue its descent toward an even lower minimum.
If the tunneling rate is sufficiently small, then the volume of space that exits inflation in this manner is replaced by the expansion of neighboring regions that do not.  The volume of inflating space never decreases, even as its fraction of the total volume may tend toward zero.



\subsection{Implications for Cosmology}



\subsection{But is it generic?} 

Eternal inflation is assumed by many in the field to be a generic feature of inflation \cite{guth07}.  What is the basis for this assumption, and how did it come to be accepted?  Does the common standard for generic agree with the one adopted in this paper?  Why or why not?  Why is the justification for the prevailing view insufficient, and what more is left to be desired?  How do we claim our results go further in addressing this question?

Check which models are compared with experimental results in the Planck and WMAP papers.

In \cite{mersiniparker07}, it is argued that application of an adiabatic regularization procedure in the calculation of the amplitude of quantum fluctuations is responsible for increasing the energy cost of stochastic eternal inflation to prohibitively high levels.


